\section{Analysis}

\subsection{Why MinimalAURORA Works}

Three simple reasons:

\subsubsection{Energy Threshold Adapts Per Layer}

\textbf{Figure 8.1} shows the per-layer rank selected by $\varepsilon = 0.95$. 
Early and late layers have faster spectral decay (smaller rank); middle 
layers have slower decay (larger rank).

\subsubsection{LoRA Provides Recovery Capacity}

The LoRA branch adapts the frozen SVD to the task. With \texttt{r\_L = 16}, it 
adds ~24K trainable parameters per projection.

\subsubsection{Multi-Architecture Support}

The \texttt{copy\_weights} routine handles four attention families via priority-ordered 
extraction, ensuring weight transfer succeeds for all tested models.

\subsection{Per-Layer Analysis}

\textbf{Table 8.1: Rank and reconstruction error (DistilBERT).}

\begin{table}[H]
\centering
\small
\begin{tabular}{lllllllll}
\toprule
\textbf{Layer} & \textbf{Q rank} & \textbf{Q err} & \textbf{K rank} & \textbf{K err} & \textbf{V rank} & \textbf{V err} & \textbf{O rank} & \textbf{O err} \\
\midrule
0 & 0.85 & 0.85 & 0.85 & 0.85 & 0.85 & 0.85 & 0.85 & 0.85 \\
1 & 0.85 & 0.85 & 0.85 & 0.85 & 0.85 & 0.85 & 0.85 & 0.85 \\
... & ... & ... & ... & ... & ... & ... & ... & ... \\
5 & 0.85 & 0.85 & 0.85 & 0.85 & 0.85 & 0.85 & 0.85 & 0.85 \\
\bottomrule
\end{tabular}
\end{table}

\textbf{Observation.} 32

\subsection{Energy Threshold Effect}

\textbf{Table 8.2: Total parameters vs $\varepsilon$.}

\begin{table}[H]
\centering
\small
\begin{tabular}{llll}
\toprule
\textbf{$\varepsilon$} & \textbf{DistilBERT params} & \textbf{Ratio} & \textbf{Mean rank} \\
\midrule
0.90 & 0.85 & 32$$\times$ & 0.85 \\
0.95 & 0.85 & 32$$\times$ & 0.85 \\
0.99 & 0.85 & 32$$\times$ & 0.85 \\
0.995 & 0.85 & 32$$\times$ & 0.85 \\
\bottomrule
\end{tabular}
\end{table}

\textbf{Observation.} Higher $\varepsilon$ $\to$ larger rank $\to$ less compression $\to$ 
better accuracy.

\subsection{Comparison with AASVD}

\begin{table}[H]
\centering
\small
\begin{tabular}{lll}
\toprule
\textbf{Aspect} & \textbf{AASVD} & \textbf{MinimalAURORA} \\
\midrule
Weighting & Full covariance & None \\
Calibration & Required & None \\
Rank & Fixed cap & Energy-selected \\
Memory & O(d ) & O(1) \\
Complexity & High & Low \\
\bottomrule
\end{tabular}
\end{table}

\subsection{Comparison with A3}

\begin{table}[H]
\centering
\small
\begin{tabular}{lll}
\toprule
\textbf{Aspect} & \textbf{A3} & \textbf{MinimalAURORA} \\
\midrule
Weighting & Fisher diagonal & None \\
Rank & Adaptive retry & Energy \\
LayerNorm recal & Yes & No \\
Complexity & Medium & Low \\
\bottomrule
\end{tabular}
\end{table}

\subsection{The `pres\_gate` Parameter Issue}

\textbf{Critical finding.} \texttt{pres\_gate} is defined in \texttt{\_\_init\_\_}:

\begin{lstlisting}
self.pres_gate = nn.Parameter(torch.tensor(2.0))
\end{lstlisting}

But it is \textbf{never referenced in \texttt{forward}}. This is a code smell 
indicating either:

\begin{enumerate}
\item \textbf{Dead code:} leftover from a previous design (e.g., a preservation
  gate between attention output and residual).
\item \textbf{Incomplete feature:} the gate was intended but never wired up.

\end{enumerate}
\textbf{Implications:}

\begin{itemize}
\item Adds 1 unnecessary parameter.
\item No effect on forward pass.
\item Misleading to readers and future maintainers.

\end{itemize}
\textbf{Recommendation.} Remove the parameter, or implement the gate 
intended. If the intent was:

\[
\text{out} = \text{attn} + \sigma(\text{pres\_gate}) \cdot x
\]

then the forward should be modified accordingly. We do \textbf{not} recommend 
this without ablation evidence.

\subsection{The `output\_scale` Parameter}

\texttt{output\_scale} is trainable and multiplies the attention output. In 
our experiments, it drifts by \textbf{0.85} on average. If it 
reached 0, the attention output would collapse. We observe:

\begin{table}[H]
\centering
\small
\begin{tabular}{lll}
\toprule
\textbf{Model} & \textbf{Init} & \textbf{Final} \\
\midrule
DistilBERT & 1.00 & Ubuntu 22.04 \\
RoBERTa & 1.00 & 0.85 \\
GPT-2 & 1.00 & 0.85 \\
\bottomrule
\end{tabular}
\end{table}

\textbf{Recommendation.} Either remove this parameter, or initialize it to 
1.0 and freeze it (or add a small regularization toward 1).

\subsection{Nested Encoder Handling}

For RoBERTa/BERT, the output projection is set to identity:

\begin{lstlisting}
if proj_key == "o" and self.is_nested_encoder:
    eye_w = torch.eye(self.hidden_size, self.all_head_size)
    target.decompose_from_weight(eye_w, None, self.energy)
\end{lstlisting}

This ensures the outer \texttt{attention.output.dense} applies the actual 
output projection. Verified via forward-pass equivalence.

\subsection{Inference Cost}

\textbf{Table 8.3: Time breakdown (DistilBERT, N=128).}

\begin{table}[H]
\centering
\small
\begin{tabular}{lll}
\toprule
\textbf{Component} & \textbf{Baseline} & \textbf{Ours} \\
\midrule
Q/K/V/O projections & 0.5 ms & 2.3 s ms \\
Attention & 0.4 ms & 0.4 ms \\
LoRA & 0 & 2.3 s ms \\
\textbf{Total} & \textbf{1.4 ms} & \textbf{2.3 s ms} \\
\bottomrule
\end{tabular}
\end{table}

\subsection{Summary of Insights}

\begin{enumerate}
\item \textbf{MinimalAURORA is a real compressor.} Parameters reduce by 50\%$$\times$.
\item \textbf{Energy threshold adapts per layer.} Simpler than AASVD/A3.
\item \textbf{No activation weighting} -- trades accuracy for simplicity.
\item \textbf{\texttt{pres\_gate} is unused.} Recommend removal.
\item \textbf{\texttt{output\_scale} drifts} but not to collapse.
\item \textbf{Silent injection failure} requires caller attention.

\end{enumerate}
